\chapter{FTP Service and File Transfer}
\chapterguide{Direct IPv4 reachability between two routers and basic VRP user/configuration views.}{Enable an FTP server, configure an authorized FTP user, connect from a client router, list files, transfer files in binary mode, verify transfers, close the session, and clean up safely.}
\sourcelab{Lab 5 - FTP.pptx}

\section{The service model used in the lab}
The source presents FTP as a file-transfer mechanism between networked devices, with the FTP server providing the service point. In this exercise R1 is the server, R2 is the client, and both routers share network \cmd{10.0.12.0/24}.

\sourceimage[0.76\linewidth]{R1--R2 FTP topology from the supplied Lab 5 slide}{ftp-topology.png}

\begin{center}
\begin{tabular}{llll}
\toprule
Device & Role & Interface & Address\\\midrule
R1 & FTP server & G0/0/1 & \cmd{10.0.12.1/24}\\
R2 & FTP client & G0/0/1 & \cmd{10.0.12.2/24}\\\bottomrule
\end{tabular}
\end{center}

\section{Step 1: prove IP connectivity before touching FTP}
\begin{lstlisting}[style=dcnterminal]
[R1] interface GigabitEthernet 0/0/1
[R1-GigabitEthernet0/0/1] ip address 10.0.12.1 24
[R2] interface GigabitEthernet 0/0/1
[R2-GigabitEthernet0/0/1] ip address 10.0.12.2 24
<R1> ping 10.0.12.2
\end{lstlisting}

\begin{keyidea}
Application troubleshooting starts below the application. If R1 and R2 cannot reach each other at IP, FTP configuration cannot succeed. The lab therefore verifies the direct network first.
\end{keyidea}

\section{Step 2: enable the FTP server and define its working directory}
The slides state that the FTP service is disabled by default. R1 explicitly enables the server and sets the default directory to \cmd{flash:/}.
\begin{lstlisting}[style=dcnterminal]
[R1] ftp server enable
[R1] set default ftp-directory flash:/
\end{lstlisting}

\section{Step 3: authorize an FTP user}
User authorization is configured under AAA. The supplied example creates user \cmd{huawei}, assigns a password, enables the FTP service type, sets privilege level 15, and assigns the FTP directory.
\begin{lstlisting}[style=dcnterminal]
[R1] aaa
[R1-aaa] local-user huawei password cipher huawei123
[R1-aaa] local-user huawei service-type ftp
[R1-aaa] local-user huawei privilege level 15
[R1-aaa] local-user huawei ftp-directory flash:/
\end{lstlisting}

\section{Step 4: verify the server process}
The source uses \cmd{display ftp-server}. Its example output shows that the server is running and that the default listening port is TCP 21.

\sourceimage[0.60\linewidth]{FTP server status output shown in the source material}{ftp-server-status.png}

\begin{verifybox}
Before attempting a login, you should be able to prove three things: the IP link works, the FTP server is running, and an authorized user exists with the required service type and directory.
\end{verifybox}

\section{Step 5: log in from R2}
\begin{lstlisting}[style=dcnterminal]
<R2> ftp 10.0.12.1
User(10.0.12.1:(none)): huawei
Enter password: ********
[R2-ftp]
\end{lstlisting}

\compactsourceimage[0.58\linewidth]{Successful FTP client login from R2}{ftp-client-login.png}

The \cmd{[R2-ftp]} prompt confirms that the device is now in the FTP client session view.

\section{Step 6: inspect, download, and upload files}
The source repeatedly uses \cmd{dir} to inspect directory contents before or after a transfer. For the configuration archive it switches to binary transfer mode.
\begin{lstlisting}[style=dcnterminal]
[R2-ftp] dir
[R2-ftp] binary
[R2-ftp] get vrpcfg.zip vrpnew.zip
[R2-ftp] bye
\end{lstlisting}
The slide notes that if \cmd{vrpcfg.zip} is not present on R1, the \cmd{save} command can be used to create the configuration file used by the exercise.

To upload a file back to the server, the source uses:
\begin{lstlisting}[style=dcnterminal]
[R2-ftp] put vrpnew.zip vrpnew2.zip
\end{lstlisting}
Afterward, the server directory is checked to verify the new file exists.

\begin{conceptmap}
\textbf{dir} = inspect remote files \quad $\vert$ \quad
\textbf{get} = server $\rightarrow$ client \quad $\vert$ \quad
\textbf{put} = client $\rightarrow$ server \quad $\vert$ \quad
\textbf{bye} = close the FTP session.
\end{conceptmap}

\section{Step 7: clean up carefully}
The lab removes the created test files after the transfer exercise.
\begin{lstlisting}[style=dcnterminal]
<R1> delete flash:/vrpnew2.zip
<R2> delete flash:/vrpnew.zip
\end{lstlisting}

\begin{pitfall}
The source explicitly warns against deleting the wrong files or erasing the entire \cmd{flash:/} directory. File-management commands are destructive. Verify the exact path and filename before confirming deletion.
\end{pitfall}

\section{A practical troubleshooting order}
\begin{netsteps}
\item Ping the server IP from the client.
\item Run \cmd{display ftp-server} on R1 and confirm the service is running.
\item Check the local user's password, service type, privilege, and FTP directory.
\item Attempt login and distinguish a connectivity failure from an authentication failure.
\item Use \cmd{dir} to verify the expected file and directory before transfer.
\item After \cmd{get} or \cmd{put}, inspect the destination directory to prove the file arrived.
\end{netsteps}

\begin{quickcheck}
\begin{enumerate}
\item Why does the lab ping R2 before configuring the FTP application?
\item Which command proves that the FTP server is running?
\item What is the direction of a \cmd{get} operation?
\item Why does the lab use \cmd{dir} both before and after transfers?
\end{enumerate}
\end{quickcheck}

\begin{chapterrecap}
The FTP lab layers an application service on top of a working IP link. R1 provides the server, AAA controls authorized access, R2 logs in as a client, \cmd{dir} inspects files, \cmd{get}/\cmd{put} move files, and final directory checks prove the transfer.
\end{chapterrecap}
